\documentclass[11pt,a4paper]{article}
\usepackage[utf8]{inputenc}
\usepackage[T1]{fontenc}
\usepackage{lmodern}
\usepackage[french]{babel}
\usepackage[margin=2.3cm]{geometry}
\usepackage{amsmath,amssymb}
\usepackage{booktabs}
\usepackage{xcolor}
\usepackage{enumitem}
\usepackage{tikz}
\usetikzlibrary{decorations.pathmorphing,patterns}
\usepackage{tcolorbox}
\usepackage{hyperref}
\hypersetup{colorlinks=true,linkcolor=blue!50!black,urlcolor=blue!50!black}

\newcommand{\dd}{\mathrm{d}}
\newcommand{\kun}{k_1(1-d)}
\newtcolorbox{alamain}{colback=yellow!8,colframe=orange!70!black,title=À faire à la main,fonttitle=\bfseries}
\newtcolorbox{idee}{colback=blue!5,colframe=blue!50!black,title=L'idée à retenir,fonttitle=\bfseries}

\title{Le système à deux ressorts écrit à la manière de Cook :\\
forces nodales, efforts résultants, solution analytique}
\author{Note complémentaire, dépôt \texttt{1D\_damage\_spring\_system}}
\date{\today}

\begin{document}
\maketitle

\begin{center}
\small Cette note reprend le chapitre 2 de Cook, \emph{Finite Element Modeling for Stress
Analysis} (§2.2, méthode directe pour l'élément barre), et l'applique pas à pas à notre
système : d'abord les forces nodales $F_{ij}$ de chaque ressort pris isolément, puis
l'assemblage et les efforts résultants, puis la construction de la solution analytique. Les
grandeurs sont celles de \texttt{parameters.py}.
\end{center}

\section{La notation de Cook pour un élément barre}

\subsection{La méthode directe}
Cook considère une barre de longueur $L$, section $A$, module $E$, avec un noeud à chaque
bout (noeuds 1 et 2). Il déplace \emph{un noeud à la fois} et calcule les forces qu'il faut
\emph{appliquer aux noeuds} pour maintenir cet état de déformation, à partir de la formule
élémentaire $\delta = FL/(AE)$, soit $F = (AE/L)\,\delta$.

\begin{center}
\begin{tikzpicture}[scale=1.0]
  \draw[thick] (0,0) rectangle (3.5,0.4);
  \node[below] at (0,-0.05) {1};
  \node[below] at (3.5,-0.05) {2};
  \draw[->,thick] (-1.1,0.2) -- (0,0.2) node[midway,above] {$F_{11}$};
  \draw[->,thick] (4.6,0.2) -- (3.5,0.2) node[midway,above] {$F_{21}$};
  \draw[->] (0,0.75) -- (0.5,0.75) node[above] {$u_1$};
  \node at (1.75,-0.7) {(a) noeud 1 déplacé de $u_1$, noeud 2 bloqué};
  \begin{scope}[xshift=7.5cm]
  \draw[thick] (0,0) rectangle (3.5,0.4);
  \node[below] at (0,-0.05) {1};
  \node[below] at (3.5,-0.05) {2};
  \draw[->,thick] (0,0.2) -- (-1.1,0.2) node[midway,above] {$F_{12}$};
  \draw[->,thick] (3.5,0.2) -- (4.6,0.2) node[midway,above] {$F_{22}$};
  \draw[->] (3.5,0.75) -- (4.0,0.75) node[above] {$u_2$};
  \node at (1.75,-0.7) {(b) noeud 2 déplacé de $u_2$, noeud 1 bloqué};
  \end{scope}
\end{tikzpicture}
\end{center}

$F_{ij}$ désigne « la force au noeud $i$ associée au déplacement du noeud $j$ ». Les $F_{ij}$
sont des \emph{amplitudes} positives ; leur sens est celui des flèches.
\begin{itemize}[nosep]
  \item Cas (a) : la barre est comprimée de $u_1$. Pour maintenir cet état il faut pousser le
        noeud 1 vers la droite et retenir le noeud 2 vers la gauche :
        $F_{11} = F_{21} = \frac{AE}{L}u_1$.
  \item Cas (b) : la barre est tendue de $u_2$. Il faut retenir le noeud 1 vers la gauche et
        tirer le noeud 2 vers la droite : $F_{12} = F_{22} = \frac{AE}{L}u_2$.
\end{itemize}

\subsection{Les forces résultantes et la matrice de raideur}
Avec la convention \emph{forces et déplacements positifs dans le même sens} (vers la droite),
les forces résultantes aux noeuds quand les deux noeuds bougent en même temps sont les sommes
signées (Cook, équation 2.2-2) :
\begin{equation}
  F_1 = F_{11} - F_{12} = \frac{AE}{L}(u_1 - u_2),\qquad
  F_2 = -F_{21} + F_{22} = \frac{AE}{L}(u_2 - u_1),
\end{equation}
soit, sous forme matricielle,
\begin{equation}
  \frac{AE}{L}\begin{bmatrix}1 & -1\\ -1 & 1\end{bmatrix}
  \begin{Bmatrix}u_1\\u_2\end{Bmatrix} = \begin{Bmatrix}F_1\\F_2\end{Bmatrix},
  \qquad\text{c'est-à-dire}\quad \mathbf k\,\mathbf d = \mathbf r .
\end{equation}
Cook en tire la règle générale : \emph{une colonne de $\mathbf k$ est le vecteur des forces
nodales qu'il faut appliquer à l'élément pour maintenir un état où ce degré de liberté vaut 1
et les autres 0}. Avec $u_1 = 1$, $u_2 = 0$, le produit $\mathbf k\mathbf d$ donne la première
colonne, $\{F_{11},\,-F_{21}\}^{\mathsf T}$.

\paragraph{Deux points à ne pas confondre.}
\begin{enumerate}[nosep]
  \item « $-F_{12}$ » ne veut pas dire que $F_{12}$ est négatif ; cela veut dire que cette force,
        d'amplitude positive, agit vers la gauche.
  \item $F_1$ et $F_2$ sont les forces \emph{appliquées à la barre par l'extérieur} (les charges
        nodales $\mathbf r$), pas les forces que la barre exerce sur les noeuds. Dans le
        vocabulaire de \texttt{explication.tex}, ce sont les forces intérieures $f^{\text{int}}$ de
        l'élément : ce que les noeuds doivent recevoir pour que l'élément reste dans cet état.
        On a toujours $F_1 + F_2 = 0$ : la barre isolée est en équilibre.
\end{enumerate}

\section{Nos deux ressorts écrits élément par élément}

\subsection{Le montage et les numéros de noeuds}
\begin{center}
\begin{tikzpicture}[scale=1.1]
  \fill[pattern=north east lines] (-0.4,-0.5) rectangle (0,0.5);
  \draw[thick] (0,-0.5) -- (0,0.5);
  \draw[thick,decorate,decoration={zigzag,segment length=6pt,amplitude=3pt}] (0,0) -- (3,0);
  \node[above] at (1.5,0.25) {élément 1 : $\kun$};
  \draw[thick,fill=gray!30] (3,-0.35) rectangle (3.8,0.35);
  \node at (3.4,0) {$m_1$};
  \draw[thick,decorate,decoration={zigzag,segment length=6pt,amplitude=3pt}] (3.8,0) -- (6.8,0);
  \node[above] at (5.3,0.25) {élément 2 : $k_2$};
  \draw[thick] (6.8,-0.4) -- (6.8,0.4);
  \draw[->,thick,red] (6.8,0) -- (7.8,0) node[right] {$u_2$ imposé};
  \node[below] at (0,-0.6) {noeud 0, $u_0 = 0$};
  \node[below] at (3.4,-0.5) {noeud 1, $u_1$};
  \node[below] at (6.8,-0.5) {noeud 2, $u_2$};
\end{tikzpicture}
\end{center}
Chaque ressort est une barre de Cook dont $AE/L$ est remplacé par sa raideur : $\kun$ pour
l'élément 1 (sécante, dépendant de $u_1$ par $d$), $k_2$ pour l'élément 2. Les déplacements
nodaux sont $u_0 = 0$ (encastrement), $u_1$ (inconnu), $u_2$ (imposé).

\subsection{Élément 1, noeuds 0 et 1}
Avec les notations de Cook appliquées aux noeuds $(0,1)$ : $F_{00} = F_{10} = \kun\,u_0 = 0$ et
$F_{01} = F_{11} = \kun\,u_1$. Les forces résultantes sont
\begin{equation}
  \begin{Bmatrix}F_0^{(1)}\\F_1^{(1)}\end{Bmatrix}
  = \kun\begin{bmatrix}1 & -1\\ -1 & 1\end{bmatrix}\begin{Bmatrix}0\\u_1\end{Bmatrix}
  = \begin{Bmatrix}-\kun\,u_1\\ +\kun\,u_1\end{Bmatrix}.
\end{equation}
Lecture : pour tenir le ressort 1 allongé de $u_1$, il faut retenir le noeud 0 vers la gauche
et tirer le noeud 1 vers la droite, avec la même amplitude $\kun\,u_1$. C'est son effort
normal $N_1$.

\subsection{Élément 2, noeuds 1 et 2}
De même avec $k_2$ :
\begin{equation}
  \begin{Bmatrix}F_1^{(2)}\\F_2^{(2)}\end{Bmatrix}
  = k_2\begin{bmatrix}1 & -1\\ -1 & 1\end{bmatrix}\begin{Bmatrix}u_1\\u_2\end{Bmatrix}
  = \begin{Bmatrix}k_2(u_1-u_2)\\ k_2(u_2-u_1)\end{Bmatrix}.
\end{equation}
Pour tenir le ressort 2 allongé de $u_2 - u_1$, il faut retenir le noeud 1 vers la gauche
et tirer le noeud 2 vers la droite, avec l'amplitude $k_2(u_2-u_1)$ : son effort normal $N_2$.
Le $F_{22}$ de Cook sur cet élément est $k_2u_2$ (ce qu'il faut tirer au noeud 2 pour l'y
amener, noeud 1 bloqué), et son $F_{12}$ est la même amplitude vue du noeud 1.

\section{L'assemblage et les efforts résultants}

\subsection{Assembler : additionner à chaque noeud}
L'étape « connect elements together » de Cook consiste à additionner, en chaque noeud, les
forces résultantes de tous les éléments qui s'y rattachent, et à égaler la somme à la force
extérieure appliquée à ce noeud. Le noeud 1 reçoit l'élément 1 et l'élément 2 ; les noeuds
0 et 2 n'en reçoivent qu'un. On obtient le système à trois degrés de liberté
\begin{equation}
  \underbrace{\begin{bmatrix}
    \kun & -\kun & 0\\
    -\kun & \kun + k_2 & -k_2\\
    0 & -k_2 & k_2
  \end{bmatrix}}_{\mathbf K}
  \underbrace{\begin{Bmatrix}0\\u_1\\u_2\end{Bmatrix}}_{\mathbf D}
  = \underbrace{\begin{Bmatrix}R_0\\ 0\\ F_2\end{Bmatrix}}_{\mathbf R}.
  \label{eq:assemble}
\end{equation}
Au second membre : $R_0$ est la réaction de l'encastrement (inconnue), $0$ parce qu'aucune
force extérieure n'est appliquée au noeud 1, $F_2$ la force que la machine applique au
noeud 2 (inconnue). Chaque colonne de $\mathbf K$ somme à zéro : une structure libre ne
résiste pas à une translation d'ensemble.

\subsection{Lire les trois lignes}
\begin{itemize}[nosep]
  \item \textbf{Ligne 1} : $R_0 = -\kun\,u_1 = -N_1$. La réaction de l'encastrement est l'opposé
        de l'effort dans le ressort 1 : l'appui retient vers la gauche.
  \item \textbf{Ligne 2} : $F_1^{(1)} + F_1^{(2)} = \kun\,u_1 + k_2(u_1-u_2) = 0$. C'est l'équilibre
        du noeud 1, la seule équation qui contienne une inconnue de déplacement. On la note
        \begin{equation}
          R(u_1;u_2) \equiv \bigl(\kun + k_2\bigr)\,u_1 - k_2\,u_2 = N_1 - N_2 .
          \label{eq:residu}
        \end{equation}
        En termes de Cook, le résidu est la somme des charges nodales que les deux éléments
        \emph{réclament} au noeud 1, somme qui doit être nulle puisque personne ne les fournit.
        Le terme $k_2u_2$ est la colonne de $\mathbf K$ associée au degré de liberté imposé,
        passée au second membre : c'est le $f^{\text{ext}}$ de \texttt{explication.tex}.
  \item \textbf{Ligne 3} : $F_2 = k_2(u_2-u_1) = N_2$. La force appliquée par la machine est
        l'effort dans le ressort 2. Elle se calcule \emph{après} avoir trouvé $u_1$, et c'est
        l'ordonnée de la courbe force--déplacement.
\end{itemize}

\subsection{Les efforts résultants}
À l'équilibre, la ligne 2 impose $N_1 = N_2$. Il n'y a qu'un seul effort dans toute la
chaîne, qu'on note $F$, et qui porte plusieurs noms selon l'endroit où on le regarde :
\begin{center}
\begin{tabular}{@{}lll@{}}
\toprule
Point de vue & Expression & Nom\\
\midrule
effort dans le ressort 1 & $\kun\,u_1$ & $N_1$, ou $F_1(u_1)$, courbe du ressort endommageable\\
effort dans le ressort 2 & $k_2(u_2-u_1)$ & $N_2$, ce que calcule le code (\texttt{F\_sol})\\
force de la machine au noeud 2 & ligne 3 de \eqref{eq:assemble} & $F_2$, la réaction, mesurée par une cellule de force\\
réaction de l'encastrement & ligne 1 de \eqref{eq:assemble} & $R_0 = -F$\\
force de la structure & $F(u_2)$ & ordonnée de la courbe force--déplacement\\
\bottomrule
\end{tabular}
\end{center}
La somme des trois lignes de \eqref{eq:assemble} donne $R_0 + 0 + F_2 = 0$ : les deux seules
forces extérieures s'équilibrent, même amplitude, sens opposés.

\paragraph{Hors équilibre.} Pendant les itérations de Newton, ou en dynamique, la ligne 2
n'est pas nulle. En statique elle vaut le résidu $R$, et la somme des lignes donne
$R_0 + R + F_2 = 0$ : la force non équilibrée sur le noeud 1 est l'écart entre ce que tire la
machine et ce que retient l'appui. En dynamique cet écart est $-m_1\ddot u_1$ : après la
rupture, $N_1 = 0$ alors que $N_2 = k_2(u_2-u_1)$ oscille entre $\pm F_{\max}$. C'est pourquoi
le code définit toujours la force par le ressort 2 : c'est la seule qui ait un sens hors
équilibre, et c'est la réaction au point de pilotage.

\subsection{Le cas endommageable : même forme, système non linéaire}
Le seul écart avec Cook est que $\kun$ dépend de $u_1$ par $d = 1-e^{-(u_1/u_c)^2}$. La
matrice garde la même forme et reste symétrique, mais la ligne 2 devient une équation
non linéaire en $u_1$. Pour la résoudre par Newton, on dérive le résidu \eqref{eq:residu} :
\begin{equation}
  k_t = \frac{\partial R}{\partial u_1} = \kun + u_1\frac{\dd[\kun]}{\dd u_1} + k_2
  = k_1e^{-x^2}(1-2x^2) + k_2,\qquad x = u_1/u_c .
\end{equation}
C'est la « colonne de $\mathbf K$ » du noeud 1, où la raideur sécante $\kun$ est remplacée par
la raideur tangente du ressort 1. L'assemblage de Cook donne donc directement la tangente du
Newton--Raphson en remplaçant, dans le bloc de l'élément 1, $\kun$ par $\kun + u_1\,\dd[\kun]/\dd u_1$.

\begin{alamain}
Écrire les deux matrices élémentaires $2\times2$, assembler le $3\times3$, et vérifier que
chaque colonne somme à zéro. Imposer $u_0 = 0$ et $u_2$ donné, isoler la ligne 2, retrouver
\eqref{eq:residu}. Calculer $F_2$ par la ligne 3 et $R_0$ par la ligne 1, vérifier $R_0 = -F_2$ en
utilisant la ligne 2.
\end{alamain}

\section{La solution analytique}

\subsection{Le problème}
À $u_2$ donné, la ligne 2 s'écrit
\begin{equation}
  k_1\,u_1\,e^{-u_1^2/u_c^2} + k_2\,u_1 = k_2\,u_2 .
  \label{eq:transc}
\end{equation}
C'est une équation \emph{transcendante} en $u_1$ : l'inconnue apparaît dans un facteur
linéaire et dans une exponentielle, il n'existe pas de formule fermée pour $u_1(u_2)$, et
selon $u_2$ elle a une ou trois racines. C'est le problème que Newton attaque numériquement.

\subsection{L'idée : inverser le rôle des variables}
L'équation \eqref{eq:transc} est \emph{linéaire en $u_2$}. Au lieu de chercher $u_1$ pour un
$u_2$ donné, on prend $u_1$ comme paramètre et on calcule le $u_2$ qui lui correspond :
\begin{equation}
  \boxed{\;u_2(u_1) = u_1 + \frac{k_1\,u_1\,e^{-u_1^2/u_c^2}}{k_2}
  = u_1\Bigl(1 + \frac{k_1}{k_2}e^{-u_1^2/u_c^2}\Bigr)\;}
  \label{eq:u2u1}
\end{equation}
Lecture : le déplacement imposé est la somme de l'allongement du ressort 1 ($u_1$) et de
l'allongement du ressort 2 sous le même effort ($N_1/k_2$). C'est la ligne 3 de Cook lue à
l'envers : $u_2 = u_1 + F_2/k_2$ avec $F_2 = N_1$.

La force le long de la branche est l'effort commun, qu'on écrit avec l'expression du
ressort 1 parce qu'elle ne dépend que du paramètre :
\begin{equation}
  \boxed{\;F(u_1) = k_1\,u_1\,e^{-u_1^2/u_c^2}\;}
  \label{eq:Fu1}
\end{equation}
On aurait pu écrire $F = k_2\bigl(u_2(u_1) - u_1\bigr)$ : en substituant \eqref{eq:u2u1} on
retrouve exactement \eqref{eq:Fu1}.

\subsection{Parcourir la branche}
On fait croître $u_1$ de $0$ vers l'infini, on calcule le couple $\bigl(u_2(u_1), F(u_1)\bigr)$
en chaque point, et on trace $F$ en fonction de $u_2$. C'est ce que fait
\texttt{analytical\_solution.analytical\_curve} : un \texttt{linspace} sur $u_1$ et les deux
expressions \eqref{eq:u2u1}--\eqref{eq:Fu1} vectorisées. Chaque $u_1$ donne un équilibre
unique, donc on obtient \emph{toute} la courbe d'équilibre, y compris la partie que ni
l'expérience ni le Newton piloté en $u_2$ ne peuvent suivre.

Le paramétrage est légitime parce que, le long de la branche, $u_1$ est strictement
croissant : c'est $u_2$ qui rebrousse chemin, pas $u_1$. L'allongement du ressort 1 ne
diminue jamais, donc la variable d'histoire vaut toujours $u_1$ et la loi en charge
$d = \hat d(u_1)$ s'applique partout.

\subsection{Lire la courbe}
La dérivée du paramétrage,
\begin{equation}
  \frac{\dd u_2}{\dd u_1} = 1 + \frac{1}{k_2}\frac{\dd F}{\dd u_1}
  = 1 + \frac{k_1}{k_2}e^{-x^2}(1-2x^2) = \frac{k_t}{k_2},
\end{equation}
dit où $u_2$ monte ($k_t > 0$) et où il redescend ($k_t < 0$).
\begin{itemize}[nosep]
  \item \textbf{Pic de force} : $\dd F/\dd u_1 = 0$, soit $x = 1/\sqrt2$, $F_{\max} = k_1u_c/\sqrt{2e}$,
        et $u_2^{\text{pic}} = u_1^{\text{pic}} + F_{\max}/k_2$.
  \item \textbf{Point limite} : première racine de $k_t = 0$, c'est-à-dire de
        $e^{-x^2}(2x^2-1) = k_2/k_1$, à chercher par dichotomie sur $[1/\sqrt2,\ \sqrt{3/2}]$.
        Au-delà, $\dd u_2/\dd u_1 < 0$ : suivre la branche exigerait de \emph{diminuer} $u_2$
        tout en allongeant le béton. C'est le snap-back.
  \item \textbf{Rebroussement bas} : seconde racine de $k_t = 0$, au-delà de $\sqrt{3/2}$, où
        $u_2$ est minimal sur la branche.
  \item \textbf{Branche rompue} : pour $x \gg 1$, $e^{-x^2}\to0$, donc $u_2 \to u_1$ et $F \to 0$.
\end{itemize}
Une instabilité n'existe que si $k_t$ peut s'annuler, donc si $k_2 < \max(-\dd F/\dd u_1) =
2e^{-3/2}k_1$ ; sinon $\dd u_2/\dd u_1 > 0$ partout et la branche est monotone en $u_2$.

\subsection{Ce que la branche permet de calculer}
\begin{itemize}[nosep]
  \item Le saut quasi-statique : à $u_2 = u_2^{\lim}$, on cherche sur la branche le second
        point de même $u_2$, loin sur la queue rompue ; le segment vertical qui les joint est
        le saut.
  \item L'énergie libérée : l'aire $|W| = \bigl|\int F\,\dd u_2\bigr|$ entre le saut vertical et
        la branche, intégrée par trapèzes le long du paramètre $u_1$. C'est ce que fait
        \texttt{released\_energy} : elle repère les deux indices de la branche ayant le même
        $u_2$ que le pic, puis intègre entre eux.
\end{itemize}

\begin{idee}
Cook écrit chaque élément, assemble, et isole les lignes des inconnues. Sur notre système
cela donne une seule équation non linéaire, la ligne du noeud 1, et deux lignes de
post-traitement pour les réactions. La solution analytique naît de la remarque que cette
équation, impossible à inverser en $u_1$, est linéaire en $u_2$ : on la lit à l'envers, et la
ligne 3 de Cook, $u_2 = u_1 + F/k_2$, devient le paramétrage de toute la branche d'équilibre.
\end{idee}

\section*{Références}
\begin{itemize}[nosep]
  \item R.~D. Cook, \emph{Finite Element Modeling for Stress Analysis}, chapitre 2, §2.2
        « Stiffness matrix formulation: bar element », méthode directe, figure 2.2-1,
        équations 2.2-1 à 2.2-3.
  \item Planches \emph{Finite Element Method, imposed displacement and applied force}
        (Varsovie) : passage de la colonne du degré de liberté imposé au second membre.
  \item Z.~P. Bažant, L. Cedolin, \emph{Stability of Structures}, figure 13.6c-d-e :
        construction de la réponse globale en ajoutant les allongements à force égale.
\end{itemize}

\end{document}
